\documentclass[11pt,a4paper]{article}

\usepackage[utf8]{inputenc}
\usepackage[T1]{fontenc}
\usepackage[brazil]{babel}
\usepackage{lmodern}
\usepackage{geometry}
\usepackage{microtype}
\usepackage{xcolor}
\usepackage{graphicx}
\usepackage{tabularx}
\usepackage{booktabs}
\usepackage{enumitem}
\usepackage{tikz}
\usepackage{fancyhdr}
\usepackage{hyperref}

\usetikzlibrary{arrows.meta,positioning}
\geometry{margin=2.3cm, headheight=14pt}
\setlength{\parindent}{0pt}
\setlength{\parskip}{0.6em}
\setlist[itemize]{leftmargin=1.5em,itemsep=0.2em}
\setlist[enumerate]{leftmargin=1.5em,itemsep=0.2em}

% ---------------------------------------------------------------
% Dados do projeto: altere aqui e o resto do documento acompanha
% ---------------------------------------------------------------
\newcommand{\jogo}{The Last Contract}
\newcommand{\subtitulo}{Mini Game Design Document}
\newcommand{\equipe}{André Bartzen Araujo \and Jackson Matiazzi \and José Maciel Camargo}
\newcommand{\disciplina}{Disciplina --- Avaliação Grau 2}
\newcommand{\gddversao}{0.23}

\definecolor{Destaque}{HTML}{7A2E2E}
\definecolor{Pendente}{HTML}{9A9A9A}

% Texto ainda não preenchido aparece em cinza e itálico.
% Para achar o que falta: buscar por \preencher no .tex
\newcommand{\preencher}[1]{\textcolor{Pendente}{\textit{[#1]}}}

\pagestyle{fancy}
\fancyhf{}
\fancyhead[L]{\textcolor{Destaque}{\jogo}}
\fancyhead[R]{\subtitulo{} --- v\gddversao}
\fancyfoot[C]{\thepage}
\renewcommand{\headrulewidth}{0.4pt}

\hypersetup{
    colorlinks=true,
    linkcolor=black,
    urlcolor=Destaque,
    pdftitle={\jogo{} - \subtitulo}
}

\title{\textbf{\jogo} \\ \large \subtitulo{} --- versão \gddversao}
\author{\equipe}
\date{\disciplina \\ \today}

\begin{document}

\maketitle
\thispagestyle{empty}
\tableofcontents
\newpage

% ===============================================================
\section{Ficha técnica}

\begin{tabularx}{\textwidth}{@{}lX@{}}
\toprule
\textbf{Título}          & \jogo \\
\textbf{Gênero}          & Ação 3D \\
\textbf{Perspectiva}     & Terceira pessoa \\
\textbf{Ambientação}     & Vila medieval tomada por mortos-vivos \\
\textbf{Estrutura}       & Mapa único, uma vida por tentativa \\
\textbf{Plataforma}      & \preencher{PC / Web} \\
\textbf{Engine}          & Godot 4.7 (GDScript) \\
\textbf{Público-alvo}    & Jogadores experientes, 16+, acostumados a ação e a escolhas com custo e benefício \\
\bottomrule
\end{tabularx}

% ===============================================================
\section{Conceito}

Jogo de ação 3D em terceira pessoa ambientado em uma vila medieval tomada por mortos-vivos. O jogador escolhe entre quatro aventureiros e recebe de um Lorde a missão de eliminar os inimigos e derrotar o boss responsável pela infestação.

A proposta é permitir que o jogador explore livremente a vila, enfrente grupos de inimigos e tente concluir a missão utilizando apenas uma vida por tentativa.

% ===============================================================
\section{História}

Uma vila medieval foi completamente dominada por mortos-vivos. Diante da ameaça, um Lorde convoca aventureiros para eliminar a infestação e derrotar a criatura responsável pelo ataque.

Caso o aventureiro morra, ele pode aceitar uma nova tentativa da missão. Se decidir desistir definitivamente, o Lorde considera o abandono uma falha grave e ordena sua execução.

% ===============================================================
\section{Gameplay}

O jogador explora livremente a vila e enfrenta grupos de mortos-vivos espalhados pelo cenário.

Cada tentativa começa com o jogador e os grupos de inimigos posicionados aleatoriamente pelo mapa. O jogador possui apenas uma vida.

Ao morrer, pode escolher entre \textbf{Reviver} ou \textbf{Desistir}. Ao reviver, todos os inimigos retornam, novas posições são sorteadas e o jogador recebe uma escolha entre dois buffs aleatórios.

Não existe sistema de níveis para jogadores ou inimigos.

\subsection{Ciclo de jogo}

\begin{center}
\begin{tikzpicture}[
    node distance=0.9cm and 1.2cm,
    etapa/.style={draw=Destaque, rounded corners, thick, minimum width=3.4cm,
                  minimum height=0.9cm, align=center, font=\small},
    fim/.style={etapa, fill=Destaque!12},
    seta/.style={-{Stealth[length=2.5mm]}, thick, Destaque},
    rotulo/.style={font=\footnotesize, midway}
]
    \node[etapa] (contrato) {Aceitar o contrato\\do Lorde};
    \node[etapa, below=of contrato] (sorteio) {Sortear posições do\\jogador e dos inimigos};
    \node[etapa, below=of sorteio] (combate) {Explorar a vila e\\enfrentar os grupos};
    \node[etapa, below=of combate] (boss) {Derrotar o boss};
    \node[fim, below=of boss] (vitoria) {Vitória};

    \node[etapa, right=2.2cm of combate] (morte) {Morte};
    \node[etapa, above=of morte] (buff) {Escolher 1 entre\\2 buffs aleatórios};
    \node[fim, below=of morte] (execucao) {Desistir:\\execução};

    \draw[seta] (contrato) -- (sorteio);
    \draw[seta] (sorteio) -- (combate);
    \draw[seta] (combate) -- (boss);
    \draw[seta] (boss) -- (vitoria);
    \draw[seta] (combate) -- (morte);
    \draw[seta] (boss.east) -- (morte.south west);
    \draw[seta] (morte) -- node[rotulo, right] {Reviver} (buff);
    \draw[seta] (buff) |- (sorteio);
    \draw[seta] (morte) -- (execucao);
\end{tikzpicture}
\end{center}

\subsection{Mecânicas}

\begin{tabularx}{\textwidth}{@{}lX@{}}
\toprule
\textbf{Mecânica} & \textbf{Descrição} \\
\midrule
Movimentação & Terceira pessoa, com exploração livre da vila. Resposta rápida: atinge a velocidade de andar (4\,m/s) em cerca de 0,1\,s e para em cerca de 0,1\,s; correr chega a 7,5\,m/s. O pulo sobe e cai num arco firme (a queda é mais rápida que a subida) e perdoa pequenos erros de tempo: pular logo após sair de uma borda ou um instante antes de aterrissar ainda funciona. \\
Combate & Corpo a corpo e à distância, conforme o aventureiro escolhido. Cada golpe tira 1 ponto de vida; o alvo reage ao dano e, ao chegar a zero, morre. \\
Vida & Pontos inteiros: aventureiros têm 10 e mortos-vivos têm 5. Uma barra de vida flutua sobre cada personagem (verde para aventureiros, vermelha para mortos-vivos). Ao levar dano, um rastro claro mostra o pedaço perdido antes de encolher e um número de dano sobe acima da cabeça. \\
Vida única & O jogador possui apenas uma vida por tentativa. \\
Grupos de inimigos & Mortos-vivos organizados em grupos espalhados pelo mapa. \\
Posições aleatórias & A cada tentativa, o jogador renasce em um ponto sorteado e os grupos de inimigos são reposicionados. \\
Reviver ou Desistir & Escolha apresentada ao morrer. Desistir encerra a missão com a execução do aventureiro. \\
Buffs acumulativos & Ao reviver, o jogador escolhe um entre dois buffs aleatórios; os buffs se acumulam entre tentativas. \\
Buffs dos inimigos & Os mortos-vivos recebem os mesmos buffs escolhidos pelo jogador nas tentativas anteriores. \\
Boss & O Necromante, um mago esqueleto que surge na praça quando todos os grupos caem: mantém distância, dispara rajadas mágicas e levanta mortos-vivos ao seu redor. \\
Objetos interativos & Baús espalhados pela vila: perto de um, aparece o aviso para abrir; aberto, solta uma poção que cura ao ser tocada. \\
Som 3D & Passos, golpes, dor, ossos caindo e feitiços saem da posição de quem os faz; a música muda conforme a cena. \\
Vitória & Limpar todos os grupos (a HUD mostra quantos faltam), derrotar o Necromante e ver a tela ``Contrato cumprido''. \\
\bottomrule
\end{tabularx}

\subsection{Controles}

O jogador escolhe em Configurações $\rightarrow$ Controles se joga com \textbf{teclado e mouse}, com \textbf{controle} (layout Xbox; testado com um 8BitDo Ultimate 2) ou no modo \textbf{automático}, o padrão, em que vale o último dispositivo usado. Com um dispositivo fixo, as ações de jogo respondem só a ele; os menus aceitam os dois sempre.

\begin{tabularx}{\textwidth}{@{}lXX@{}}
\toprule
\textbf{Ação} & \textbf{Teclado e mouse} & \textbf{Controle} \\
\midrule
Mover      & W, A, S, D (relativo à câmera) & Analógico esquerdo \\
Correr     & W duas vezes seguidas; corre enquanto W estiver pressionado & Clicar no analógico esquerdo; corre até soltar o analógico \\
Pular      & Espaço (vale apertar um pouco antes de tocar o chão) & A \\
Câmera     & Movimento do mouse (sobre o ombro direito) & Analógico direito \\
Atacar     & Botão esquerdo; segurar ataca em sequência e um clique pouco antes da recarga terminar fica guardado & RT (mesmas regras) \\
Habilidade & Botão direito: o Knight ergue o escudo enquanto segura; o Barbarian golpeia com o segundo machado & LT \\
Interagir  & E (abrir baú) & X \\
Pausar     & Esc & Menu (Start) \\
Menus      & Mouse, setas, Enter, Esc; teclas 1--4 na seleção e 1--2 nas bênçãos & Direcional, A confirma, B volta, LB/RB trocam aventureiro ou aba \\
\bottomrule
\end{tabularx}

As ligações dos dois dispositivos ficam no Input Map do projeto e podem ser trocadas no editor. Sensibilidade do mouse e do analógico e a inversão do eixo vertical também ficam nas configurações. A cena \texttt{scenes/tools/gamepad\_tester.tscn} mostra ao vivo os botões, analógicos, gatilhos e ações recebidos, para conferir um controle novo.

\textbf{Problema conhecido:} um 8BitDo Ultimate 2 ligado por Bluetooth aparece como conectado, mas não manda nenhum comando ao Godot 4.7, embora funcione no painel de controles do Windows. A causa é o driver de 8BitDo da biblioteca SDL, que o Godot usa para ler controles. A solução é definir a variável de ambiente \texttt{SDL\_JOYSTICK\_HIDAPI\_8BITDO=0} (\texttt{setx SDL\_JOYSTICK\_HIDAPI\_8BITDO 0} e reabrir a Steam e o Godot) ou ligar o controle pelo receptor 2.4G ou por cabo.

\subsection{Buffs}

Ao reviver, o jogador escolhe uma entre duas bênçãos sorteadas. Pelo contrato do Lorde, a bênção escolhida vale para o aventureiro \emph{e também para todos os mortos-vivos} nas tentativas seguintes. Como o aventureiro tem 10 de vida e enfrenta vários mortos-vivos de 5 de vida, o mesmo buff pesa diferente para cada lado: a escolha é decidir qual efeito favorece mais o jogador do que os inimigos. Metade dos buffs tem custo e benefício, pensados para o público experiente.

\begin{tabularx}{\textwidth}{@{}lX@{}}
\toprule
\textbf{Buff} & \textbf{Efeito (em todos os personagens)} \\
\midrule
Lâmina Afiada & +1 de dano \\
Couraça & +3 de vida máxima \\
Passos Leves & +20\% de velocidade \\
Frenesi & $-$25\% de tempo entre ataques, $-$2 de vida máxima \\
Fúria Cega & +1 de dano, $-$3 de vida máxima \\
Sangue Frio & +2 de vida máxima, $-$10\% de velocidade \\
\bottomrule
\end{tabularx}

Regras: cada buff acumula até 3 vezes; os dois sorteados são sempre diferentes; vida máxima nunca fica abaixo de 1; velocidade e tempo entre ataques nunca caem abaixo de 20\% do valor base. Desistir na tela de morte, ou abandonar o contrato pela pausa, leva à cena de execução do aventureiro e depois ao menu principal; o próximo contrato começa do zero (tentativas e buffs zerados).

% ===============================================================
\section{Personagens}

O jogador poderá escolher entre quatro aventureiros, com diferentes estilos de combate, incluindo personagens corpo a corpo e de longo alcance.

\begin{tabularx}{\textwidth}{@{}lX@{}}
\toprule
\textbf{Personagem} & \textbf{Papel} \\
\midrule
Aventureiros & Personagens controlados pelos jogadores. \\
Lorde & Responsável por entregar a missão. \\
Mortos-vivos & Inimigos espalhados pela vila. \\
Boss & Principal inimigo e responsável pela infestação. \\
\bottomrule
\end{tabularx}

\subsection{Aventureiros jogáveis}

Cada aventureiro tem um ataque único. Todos atacam na direção em que a câmera olha. Mage e Rogue têm uma mira no centro da tela: o projétil sai da arma em direção ao ponto que a mira cobre (parede, chão ou morto-vivo, até 40\,m), inclusive para cima ou para baixo, então acertar de longe exige apontar a câmera; todos causam 1 de dano por acerto (mais os buffs) e nenhum acerta outro aventureiro. Projéteis param ao bater em paredes e deixam rastro: o projétil do Mage queima uma marca violeta na superfície (some após 6\,s) e o virote do Rogue fica cravado onde bateu, na parede ou no corpo do morto-vivo, acompanhando-o enquanto ele se move (some após 4\,s). Golpes corpo a corpo não atravessam paredes: só acertam quem está com o caminho livre até o atacante. O virote da besta fica cravado no osso mais próximo de onde acertou e acompanha o morto-vivo andando, girando e caindo ao morrer.

Regras comuns de ataque, pensadas para o controle não ficar pesado:
\begin{itemize}
    \item durante o ataque o aventureiro ainda se move a 30\% da velocidade, sem girar, para manter a mira;
    \item logo depois do impacto, mover, pular ou atacar de novo corta o resto da animação; parado, a animação termina normalmente;
    \item ser atingido durante a preparação (antes do impacto) cancela o golpe; depois do impacto, o golpe já contou;
    \item ao ser atingido, o aventureiro recupera o controle em 0,3\,s.
\end{itemize}


\begin{tabularx}{\textwidth}{@{}llX@{}}
\toprule
\textbf{Aventureiro} & \textbf{Alcance} & \textbf{Estilo de combate} \\
\midrule
Knight & Corpo a corpo & Espada e escudo. \textbf{Combo de três golpes} (vertical, diagonal, horizontal) que segue enquanto o jogador continua atacando e reinicia após uma pausa; acerta só à frente; recarga de 0,6\,s. \textbf{Escudo} (botão direito): enquanto segura, bloqueia golpes e projéteis que chegam pela frente (arco de 140°), anda devagar e não ataca; o bloqueio aparece como ``Bloqueado'' sobre a barra de vida \\
Barbarian & Corpo a corpo & Dois machados. \textbf{Giro}: acerta todos os mortos-vivos em volta; recarga de 1,6\,s. \textbf{Segundo machado} (botão direito): golpe rápido só à frente, recarga de 0,5\,s, para finalizar um alvo sem esperar o giro \\
Mage & Longo alcance & Cajado. \textbf{Projétil mágico}: mais lento (14\,m/s), alcance de 18\,m e \emph{atravessa} os inimigos em linha; recarga de 1,0\,s \\
Rogue & Longo alcance & Besta. \textbf{Virote}: rápido (30\,m/s), alcance de 24\,m, para no primeiro inimigo; recarga de 0,7\,s \\
\bottomrule
\end{tabularx}

\subsection{Mortos-vivos e IA}

Os mortos-vivos não se limitam a andar e bater: eles percebem, reagem com atraso humano e, nas próximas fases, se coordenam em grupo e respondem às ações do jogador. Todos usam a mesma base de física e animação do jogador e são mais lentos que ele (andar 2,8\,m/s, correr 5\,m/s), para que fugir seja uma opção real.

\textbf{Percepção.} Visão em cone de 120\textdegree{} até 14\,m, bloqueada por paredes; audição de corrida e de ataques num raio de 6\,m, mesmo sem ver; percepção de proximidade: nota o jogador a até 2,5\,m mesmo de costas, se não houver parede; ser atingido, inclusive por um disparo de longe, revela a posição do jogador; memória da última posição do jogador por 3\,s.

\textbf{Comportamento (implementado).}
\begin{itemize}
    \item \textbf{Adormecido}: deitado no chão; não enxerga, mas desperta com o jogador a 7\,m, com barulho ou ao ser atingido. Enquanto dorme não mostra barra de vida; ela aparece quando a animação de despertar avisa que o perigo começou. Atingido deitado, começa a levantar na hora. Deitado ou levantando, leva o dano normalmente mas não se encolhe com o golpe, então a animação de levantar nunca é cortada.
    \item \textbf{Em guarda}: ao perceber o jogador, espera 0,3\,s antes de agir, para não parecer onisciente.
    \item \textbf{Perseguição}: segue o jogador pela navegação, contornando paredes; corre quando está longe.
    \item \textbf{Ataque}: só golpeia quase de frente para o jogador, com aviso de 0,6\,s entre o início do golpe e o impacto (tempo para esquivar ou cortar o golpe); recarga de 1,8\,s.
    \item \textbf{Busca}: ao perder o jogador, vai até a última posição conhecida, procura por 2,5\,s e volta à guarda.
    \item \textbf{Corpo}: o morto-vivo derrubado fica caído por 6\,s, afunda no chão e sai da cena, deixando só uma mancha escura onde caiu. Mantém a vila limpa e reduz o que precisa ser desenhado.
    \item \textbf{Elite}: mortos-vivos podem formar um grupo. Quando o último do grupo cai, os corpos ficam parados por 2\,s, deslizam até o ponto do grupo, afundam juntos e dali se levanta um esqueleto elite: maior, com elmo de chifres e aura roxa, 15 de vida, golpe de 2 de dano, visão e audição maiores. Ele já levanta sabendo onde o jogador está e vai atrás dele. Limpar uma área tem um último confronto, e o elite serve de base para o boss. Para montar um grupo na fase, basta colocar os mortos-vivos como filhos de um nó \texttt{EnemyGroup}: cada um só pode ter um pai, então grupos não se misturam. O elite nasce no ponto onde o nó do grupo foi colocado; só os corpos a até 8\,m desse ponto são puxados (quem caiu longe some como um corpo comum), e o editor mostra um aviso se algum membro for posicionado longe demais.
\end{itemize}

\textbf{Boss: o Necromante.} Um esqueleto mago maior, com cajado e aura verde, 30 de vida. Surge na praça, levantando do chão, quando o último grupo da vila cai. Usa o mesmo cérebro dos outros mortos-vivos, mas com alcance de 11\,m: fica longe e dispara rajadas mágicas (2 de dano, recarga de 2\,s). A cada 9\,s de caça levanta dois esqueletos ao seu redor, com no máximo quatro vivos; quando ele morre, os invocados desabam junto. A HUD mostra uma barra de vida grande com o nome dele.

\textbf{Próximas fases.}
\begin{enumerate}
    \item \textbf{Grupo}: no máximo alguns mortos-vivos atacam ao mesmo tempo, os demais cercam o jogador por lados diferentes; quem vê o jogador alerta os vizinhos.
    \item \textbf{Reações}: ler o golpe do jogador; bloqueio frontal (Warrior com escudo); esquiva; atacar quando o jogador erra um golpe. As chances de reação crescem 5\% por tentativa, até 60\%.
    \item \textbf{Variantes}: Warrior (escudo), Rogue (besta que mira prevendo o movimento e mantém distância) e Mage (projétil), cada um com seu perfil de números.
\end{enumerate}

% ===============================================================
\section{Mundo}

O jogo acontece em uma vila medieval explorável, composta por ruas, casas, praça, construções e áreas abertas.

Não existem fases separadas. Toda a missão acontece em um único mapa.

A cada nova tentativa, tanto o jogador quanto os grupos de inimigos recebem novas posições aleatórias dentro da vila, fazendo com que os encontros e o caminho percorrido possam mudar entre as partidas.

\textbf{A vila (primeira versão).} Um vilarejo medieval de 48\,m $\times$ 48\,m, à noite, iluminado por tochas nas paredes. No centro, uma praça calçada com carroças, caixotes e barris que servem de cobertura; em volta, oito casas de pedra e enxaimel, de um ou dois andares, formando ruas e becos; ao sul, um campo aberto cercado. As portas em arco são vãos de verdade: dá para entrar nas casas, que viram esconderijos e atalhos.

\textbf{Sorteio dos encontros.} O mapa tem seis pontos de nascimento do jogador e treze pontos de grupo. A cada tentativa o jogador nasce num ponto sorteado e cinco grupos ocupam pontos a pelo menos 16\,m dele. Os grupos são cenas prontas: \textbf{patrulha} (2 esqueletos, sem elite), \textbf{guarda} (3 esqueletos, com elite) e \textbf{dormentes} (3 esqueletos deitados, com elite). Um teste garante que de todo ponto de grupo se anda até todo ponto de nascimento.

% Para incluir um mapa ou esboço, salve a imagem em docs/img/ e ajuste o nome.
\IfFileExists{img/mapa.png}{%
\begin{figure}[h]
    \centering
    \includegraphics[width=0.8\textwidth]{img/mapa.png}
    \caption{Esboço da vila}
\end{figure}}{}

% ===============================================================
\section{Interface}

Todas as telas compartilham um tema único de madeira escura e dourado. Ações destrutivas usam borda vermelha e sempre pedem confirmação; a ação principal de cada tela vem destacada e recebe o foco do teclado. A interface foi desenhada para 1152$\times$648 e escala com o tamanho da janela, mantendo a proporção.

\textbf{Fluxo de telas:} Menu principal $\rightarrow$ Seleção de aventureiro $\rightarrow$ Entrada na missão $\rightarrow$ Partida (HUD, Pausa) $\rightarrow$ Morte $\rightarrow$ Reviver (nova tentativa na mesma partida) ou Desistir $\rightarrow$ Execução $\rightarrow$ Menu principal. Abandonar o contrato pela pausa também leva à Execução.

\begin{tabularx}{\textwidth}{@{}lX@{}}
\toprule
\textbf{Tela} & \textbf{Conteúdo} \\
\midrule
Menu principal & Título, ``Aceitar contrato'', Configurações, Sair e créditos, sobre uma cena 3D da vila com câmera girando \\
Seleção de aventureiro & Os quatro aventureiros lado a lado em 3D; o escolhido dá um passo à frente, ganha um anel dourado no chão e o nome destacado. Ficha com estilo, ataque e descrição. Escolha por clique, teclas 1 a 4 ou setas; a escolha vale para todas as tentativas do contrato \\
Configurações & Áudio (geral, música, efeitos), vídeo (tela cheia, VSync) e controles (sensibilidade do mouse, inverter eixo vertical, lista de teclas). Salvas automaticamente \\
HUD & Tentativa atual, vida do jogador em número e barra, e bênçãos ativas com efeito e acúmulos. Para Mage e Rogue, uma mira no centro da tela indica para onde o projétil vai \\
Barras 3D & Sobre os mortos-vivos (vermelha, com rastro do dano); escondida enquanto o morto-vivo dorme, aparece quando ele começa a levantar. O jogador usa a barra do HUD. Números de dano sobem de todos os personagens \\
Entrada na missão & Cena curta à noite: o aventureiro escolhido chega pela estrada até a entrada da vila (``Contrato aceito.'', ``O Lorde quer a vila livre dos mortos.''); depois começa a partida. Qualquer tecla pula \\
Pausa & Continuar, Configurações e Abandonar contrato (com confirmação; leva à execução) \\
Morte & Tentativa, bênçãos acumuladas, Reviver ou Desistir (com confirmação; leva à execução). Reviver mostra duas cartas de bênção, escolhidas com o mouse ou as teclas 1 e 2 \\
Execução & Cena em duas tomadas: o aventureiro escolhido vai embora pela estrada da vila ao entardecer (``Ele deu as costas ao contrato''); corta para a praça pública, onde ele está ajoelhado num cadafalso diante da multidão, ouve o veredito do Lorde e recebe o golpe de machado do carrasco encapuzado, enquanto a multidão comemora; termina em ``Contrato encerrado'' e volta ao menu. Qualquer tecla pula \\
\bottomrule
\end{tabularx}

\subsection{Áudio}

Todo som de personagem é posicional (3D): passos contados pela distância andada no chão, golpe saindo no momento do impacto, dor ao levar dano, estalo no escudo ao bloquear, ossos caindo na morte do esqueleto, rangido ao levantar do chão e um sino grave quando o elite se forma. Virotes e feitiços fazem barulho onde batem. Cada grupo de sons sorteia entre variações e muda levemente o tom, para não repetir. Os botões de todos os menus têm som de clique e de foco. A música segue a cena: uma melodia medieval no menu e na seleção, um ambiente sombrio na vila e nas cutscenes. Música e efeitos têm volumes próprios nas configurações.

% ===============================================================
\section{Arquitetura técnica}

\subsection{Personagens}

Todos os personagens compartilham uma classe base. Cada classe filha só decide a \emph{intenção} de movimento; física, vida e animação ficam na base.

\begin{center}
\begin{tikzpicture}[
    node distance=0.8cm and 0.6cm,
    classe/.style={draw=Destaque, rounded corners, thick, align=center, font=\small, inner sep=5pt},
    seta/.style={-{Stealth[length=2.5mm]}, thick, Destaque}
]
    \node[classe] (character) {\textbf{Character}\\ física, vida, dano, morte, animação};
    \node[classe, below left=of character] (adventurer) {\textbf{Adventurer}\\ teclado, mouse, câmera};
    \node[classe, below right=of character] (enemy) {\textbf{Enemy}\\ IA \preencher{a implementar}};
    \node[classe, below=of enemy] (boss) {\textbf{Boss} \preencher{a implementar}};
    \draw[seta] (adventurer) -- (character);
    \draw[seta] (enemy) -- (character);
    \draw[seta] (boss) -- (enemy);
\end{tikzpicture}
\end{center}

Cada aventureiro e cada tipo de morto-vivo é uma cena herdada da base, que troca apenas o modelo, as armas visíveis e os atributos. A IA dos mortos-vivos é dividida em percepção (o que sabem), cérebro com máquina de estados (o que decidem, a cada 0,15\,s) e navegação (como se movem); os números de cada tipo ficam num perfil de dados. O ataque é um componente separado, reaproveitado por aventureiros e inimigos: \texttt{Attack} (recarga, animação e momento do impacto) com as variantes \texttt{MeleeAttack} (o formato da hitbox define golpe frontal ou giro) e \texttt{RangedAttack} (dispara um \texttt{Projectile}). As telas apenas emitem sinais; o \texttt{GameSession} decide trocas de cena e reinícios. Buffs são recursos de dados (\texttt{Buff} com uma lista de \texttt{StatModifier}); o estado da missão (tentativa e buffs acumulados) fica no \texttt{RunState}, que sobrevive à recarga da cena a cada Reviver. Quem pode ser atingido é definido pelas camadas de física: mundo, aventureiros e inimigos.

\subsection{Mapa}

A vila é montada com dois GridMaps: \textbf{Chão} (células de 2\,m para ruas e praça) e \textbf{Estrutura} (células de 1\,m para pisos, paredes, cantos e escadas das casas). Telhados, portas, janelas e objetos de cenário ficam fora da grade, como cenas avulsas. As bibliotecas de peças são geradas por um script a partir dos modelos importados. Cada mapa é uma cena própria em \texttt{scenes/maps/} (\texttt{village}, \texttt{grid\_sandbox}, \texttt{public\_square} e \texttt{village\_road}) e é desenhado inteiro pela paleta do GridMap: cada grade é uma cena separada em \texttt{scenes/maps/<mapa>/} (\texttt{ground}, \texttt{structure} e \texttt{objects}), que pode ser reaproveitada em outras cenas. São três bibliotecas de peças: \textbf{ground} (pisos), \textbf{village\_structure} (paredes, cantos, escadas, telhados e empenas) e \textbf{props} (carroça, caixotes, barris, cercas e tocha de parede, com colisão). Fora da grade ficam só as luzes, os limites invisíveis do mapa, os pontos de sorteio e peças com personagens, como o cadafalso da execução.

As cutscenes usam o AnimationPlayer do próprio Godot: câmeras, legendas, fades, movimentos e animações dos personagens são trilhas da animação, ajustáveis no painel Animation sem mexer em código. Todas usam o mesmo script (\texttt{AdventurerCutscene}), que só coloca o aventureiro escolhido nos marcadores indicados e, ao terminar, abre a cena definida no inspetor; uma cutscene nova não precisa de código.

\subsection{Testes automatizados}

O projeto usa o gdUnit4 (versão 6.2.1). São 261 testes em 36 suítes, incluindo a mira dos ataques à distância, organizados por área: seleção de aventureiro, criação do aventureiro escolhido na partida, cena de execução, um teste específico para cada personagem (Knight, Barbarian, Mage, Rogue e Skeleton Minion), personagem base (vida, morte, buffs), percepção e cérebro dos mortos-vivos (visão, audição, memória, despertar, perseguição, ataque com aviso, busca e navegação contornando paredes), habilidades do botão direito (escudo do Knight, segundo machado do Barbarian), jogo pelo controle (analógicos, gatilhos, filtro por dispositivo, navegação nos menus), formação do elite, movimento do aventureiro (aceleração, andar, correr, pulo com tolerâncias, mira pela câmera), comandos de ataque (combo, cancelamento, ataque guardado, golpe interrompido), combate corpo a corpo e à distância, estado da missão, fase da vila (sorteio dos encontros e alcance pela navegação), áudio, baú e poção, Necromante e vitória, configurações, HUD, barra de vida, tela de morte e pausa. Para rodar: \texttt{tests/run\_tests.ps1}.

% ===============================================================
\section{Assets}

O projeto utiliza pacotes gratuitos de modelos 3D e áudio, todos sob licença \href{https://creativecommons.org/publicdomain/zero/1.0/}{CC0} (domínio público, crédito opcional):

\begin{tabularx}{\textwidth}{@{}lX@{}}
\toprule
\textbf{Pacote} & \textbf{Uso no jogo} \\
\midrule
\href{https://kaylousberg.itch.io/kaykit-adventurers}{KayKit Adventurers} (Kay Lousberg) & Aventureiros jogáveis \\
\href{https://kaylousberg.itch.io/kaykit-skeletons}{KayKit Skeletons} (Kay Lousberg) & Mortos-vivos \\
\href{https://kaylousberg.itch.io/kaykit-dungeon-pack}{KayKit Dungeon Pack} (Kay Lousberg) & Props e elementos de cenário \\
\href{https://quaternius.itch.io/medieval-village-megakit}{Medieval Village MegaKit} (Quaternius) & Construção da vila \\
\href{https://kenney.nl/assets/rpg-audio}{RPG Audio}, \href{https://kenney.nl/assets/impact-sounds}{Impact Sounds} e \href{https://kenney.nl/assets/interface-sounds}{Interface Sounds} (Kenney) & Efeitos sonoros \\
\href{https://opengameart.org/content/dark-cavern-ambient}{Dark Cavern Ambient} (Paul Wortmann) & Música ambiente da vila \\
\href{https://opengameart.org/content/medieval-fair-loop}{Medieval fair [Loop]} (Woli34) & Música do menu \\
\bottomrule
\end{tabularx}

\section{Publicação}

O jogo é exportado para Windows e Linux (64 bits) como um executável único, com o conteúdo embutido, e publicado no itch.io. A versão para navegador não é usada porque o jogo depende do renderizador Forward+ (sombras, névoa e luzes dinâmicas), que o navegador não suporta. O build de Linux e este documento em PDF são gerados pelo GitHub Actions (workflows \texttt{build-linux.yml} e \texttt{gdd.yml}).

O executável usa o ícone do jogo (o Knight sob a luz verde do Necromante). O README do projeto traz os controles, como rodar, como testar e os créditos, e o workflow \texttt{.github/workflows/tests.yml} roda a suíte de testes a cada push e pull request.

\end{document}
